\IfFileExists{stacks-project.cls}{%
\documentclass{stacks-project}
}{%
\documentclass{amsart}
}

% For dealing with references we use the comment environment
\usepackage{verbatim}
\newenvironment{reference}{\comment}{\endcomment}
%\newenvironment{reference}{}{}
\newenvironment{slogan}{\comment}{\endcomment}
\newenvironment{history}{\comment}{\endcomment}

% For commutative diagrams we use Xy-pic
\usepackage[all]{xy}

% We use 2cell for 2-commutative diagrams.
\xyoption{2cell}
\UseAllTwocells

% We use multicol for the list of chapters between chapters
\usepackage{multicol}

% This is generally recommended for better output
\usepackage{lmodern}
\usepackage[T1]{fontenc}

% For cross-file-references
\usepackage{xr-hyper}

% Package for hypertext links:
\usepackage{hyperref}

% For any local file, say "hello.tex" you want to link to please
% use \externaldocument[hello-]{hello}
\externaldocument[introduction-]{introduction}
\externaldocument[conventions-]{conventions}
\externaldocument[sets-]{sets}
\externaldocument[categories-]{categories}
\externaldocument[topology-]{topology}
\externaldocument[sheaves-]{sheaves}
\externaldocument[sites-]{sites}
\externaldocument[stacks-]{stacks}
\externaldocument[fields-]{fields}
\externaldocument[algebra-]{algebra}
\externaldocument[brauer-]{brauer}
\externaldocument[homology-]{homology}
\externaldocument[derived-]{derived}
\externaldocument[simplicial-]{simplicial}
\externaldocument[more-algebra-]{more-algebra}
\externaldocument[smoothing-]{smoothing}
\externaldocument[modules-]{modules}
\externaldocument[sites-modules-]{sites-modules}
\externaldocument[injectives-]{injectives}
\externaldocument[cohomology-]{cohomology}
\externaldocument[sites-cohomology-]{sites-cohomology}
\externaldocument[dga-]{dga}
\externaldocument[dpa-]{dpa}
\externaldocument[sdga-]{sdga}
\externaldocument[hypercovering-]{hypercovering}
\externaldocument[schemes-]{schemes}
\externaldocument[constructions-]{constructions}
\externaldocument[properties-]{properties}
\externaldocument[morphisms-]{morphisms}
\externaldocument[coherent-]{coherent}
\externaldocument[divisors-]{divisors}
\externaldocument[limits-]{limits}
\externaldocument[varieties-]{varieties}
\externaldocument[topologies-]{topologies}
\externaldocument[descent-]{descent}
\externaldocument[perfect-]{perfect}
\externaldocument[more-morphisms-]{more-morphisms}
\externaldocument[flat-]{flat}
\externaldocument[groupoids-]{groupoids}
\externaldocument[more-groupoids-]{more-groupoids}
\externaldocument[etale-]{etale}
\externaldocument[chow-]{chow}
\externaldocument[intersection-]{intersection}
\externaldocument[pic-]{pic}
\externaldocument[weil-]{weil}
\externaldocument[adequate-]{adequate}
\externaldocument[dualizing-]{dualizing}
\externaldocument[duality-]{duality}
\externaldocument[discriminant-]{discriminant}
\externaldocument[derham-]{derham}
\externaldocument[local-cohomology-]{local-cohomology}
\externaldocument[algebraization-]{algebraization}
\externaldocument[curves-]{curves}
\externaldocument[resolve-]{resolve}
\externaldocument[models-]{models}
\externaldocument[functors-]{functors}
\externaldocument[equiv-]{equiv}
\externaldocument[pione-]{pione}
\externaldocument[etale-cohomology-]{etale-cohomology}
\externaldocument[proetale-]{proetale}
\externaldocument[relative-cycles-]{relative-cycles}
\externaldocument[more-etale-]{more-etale}
\externaldocument[trace-]{trace}
\externaldocument[crystalline-]{crystalline}
\externaldocument[spaces-]{spaces}
\externaldocument[spaces-properties-]{spaces-properties}
\externaldocument[spaces-morphisms-]{spaces-morphisms}
\externaldocument[decent-spaces-]{decent-spaces}
\externaldocument[spaces-cohomology-]{spaces-cohomology}
\externaldocument[spaces-limits-]{spaces-limits}
\externaldocument[spaces-divisors-]{spaces-divisors}
\externaldocument[spaces-over-fields-]{spaces-over-fields}
\externaldocument[spaces-topologies-]{spaces-topologies}
\externaldocument[spaces-descent-]{spaces-descent}
\externaldocument[spaces-perfect-]{spaces-perfect}
\externaldocument[spaces-more-morphisms-]{spaces-more-morphisms}
\externaldocument[spaces-flat-]{spaces-flat}
\externaldocument[spaces-groupoids-]{spaces-groupoids}
\externaldocument[spaces-more-groupoids-]{spaces-more-groupoids}
\externaldocument[bootstrap-]{bootstrap}
\externaldocument[spaces-pushouts-]{spaces-pushouts}
\externaldocument[spaces-chow-]{spaces-chow}
\externaldocument[groupoids-quotients-]{groupoids-quotients}
\externaldocument[spaces-more-cohomology-]{spaces-more-cohomology}
\externaldocument[spaces-simplicial-]{spaces-simplicial}
\externaldocument[spaces-duality-]{spaces-duality}
\externaldocument[formal-spaces-]{formal-spaces}
\externaldocument[restricted-]{restricted}
\externaldocument[spaces-resolve-]{spaces-resolve}
\externaldocument[formal-defos-]{formal-defos}
\externaldocument[defos-]{defos}
\externaldocument[cotangent-]{cotangent}
\externaldocument[examples-defos-]{examples-defos}
\externaldocument[algebraic-]{algebraic}
\externaldocument[examples-stacks-]{examples-stacks}
\externaldocument[stacks-sheaves-]{stacks-sheaves}
\externaldocument[criteria-]{criteria}
\externaldocument[artin-]{artin}
\externaldocument[quot-]{quot}
\externaldocument[stacks-properties-]{stacks-properties}
\externaldocument[stacks-morphisms-]{stacks-morphisms}
\externaldocument[stacks-limits-]{stacks-limits}
\externaldocument[stacks-cohomology-]{stacks-cohomology}
\externaldocument[stacks-perfect-]{stacks-perfect}
\externaldocument[stacks-introduction-]{stacks-introduction}
\externaldocument[stacks-more-morphisms-]{stacks-more-morphisms}
\externaldocument[stacks-geometry-]{stacks-geometry}
\externaldocument[moduli-]{moduli}
\externaldocument[moduli-curves-]{moduli-curves}
\externaldocument[examples-]{examples}
\externaldocument[exercises-]{exercises}
\externaldocument[guide-]{guide}
\externaldocument[desirables-]{desirables}
\externaldocument[coding-]{coding}
\externaldocument[obsolete-]{obsolete}
\externaldocument[fdl-]{fdl}
\externaldocument[index-]{index}

% Theorem environments.
%
\theoremstyle{plain}
\newtheorem{theorem}[subsection]{Theorem}
\newtheorem{proposition}[subsection]{Proposition}
\newtheorem{lemma}[subsection]{Lemma}

\theoremstyle{definition}
\newtheorem{definition}[subsection]{Definition}
\newtheorem{example}[subsection]{Example}
\newtheorem{exercise}[subsection]{Exercise}
\newtheorem{situation}[subsection]{Situation}

\theoremstyle{remark}
\newtheorem{remark}[subsection]{Remark}
\newtheorem{remarks}[subsection]{Remarks}

\numberwithin{equation}{subsection}

% Macros
%
\def\lim{\mathop{\mathrm{lim}}\nolimits}
\def\colim{\mathop{\mathrm{colim}}\nolimits}
\def\Spec{\mathop{\mathrm{Spec}}}
\def\Hom{\mathop{\mathrm{Hom}}\nolimits}
\def\Ext{\mathop{\mathrm{Ext}}\nolimits}
\def\SheafHom{\mathop{\mathcal{H}\!\mathit{om}}\nolimits}
\def\SheafExt{\mathop{\mathcal{E}\!\mathit{xt}}\nolimits}
\def\Sch{\mathit{Sch}}
\def\Mor{\mathop{\mathrm{Mor}}\nolimits}
\def\Ob{\mathop{\mathrm{Ob}}\nolimits}
\def\Sh{\mathop{\mathit{Sh}}\nolimits}
\def\NL{\mathop{N\!L}\nolimits}
\def\CH{\mathop{\mathrm{CH}}\nolimits}
\def\proetale{{pro\text{-}\acute{e}tale}}
\def\etale{{\acute{e}tale}}
\def\QCoh{\mathit{QCoh}}
\def\Ker{\mathop{\mathrm{Ker}}}
\def\Im{\mathop{\mathrm{Im}}}
\def\Coker{\mathop{\mathrm{Coker}}}
\def\Coim{\mathop{\mathrm{Coim}}}

% Boxtimes
%
\DeclareMathSymbol{\boxtimes}{\mathbin}{AMSa}{"02}

%
% Macros for moduli stacks/spaces
%
\def\QCohstack{\mathcal{QC}\!\mathit{oh}}
\def\Cohstack{\mathcal{C}\!\mathit{oh}}
\def\Spacesstack{\mathcal{S}\!\mathit{paces}}
\def\Quotfunctor{\mathrm{Quot}}
\def\Hilbfunctor{\mathrm{Hilb}}
\def\Curvesstack{\mathcal{C}\!\mathit{urves}}
\def\Polarizedstack{\mathcal{P}\!\mathit{olarized}}
\def\Complexesstack{\mathcal{C}\!\mathit{omplexes}}
% \Pic is the operator that assigns to X its picard group, usage \Pic(X)
% \Picardstack_{X/B} denotes the Picard stack of X over B
% \Picardfunctor_{X/B} denotes the Picard functor of X over B
\def\Pic{\mathop{\mathrm{Pic}}\nolimits}
\def\Picardstack{\mathcal{P}\!\mathit{ic}}
\def\Picardfunctor{\mathrm{Pic}}
\def\Deformationcategory{\mathcal{D}\!\mathit{ef}}
\usepackage{cleveref}
\crefname{originallemma}{original lemma}{original lemmas}
\Crefname{originallemma}{Original lemma}{Original lemmas}
\makeatletter
\newlabel{lemma-relative-cycles-equal}{{0H5C}{source}{}{sourcepointer-1}{}}
\newlabel{lemma-relative-cycles-equal@cref}{{[originallemma][0][]0H5C}{[1][0][]source}{}{}{}}
\makeatother
\crefname{originallemma}{original lemma}{original lemmas}
\Crefname{originallemma}{Original lemma}{Original lemmas}
\makeatletter
\newlabel{limits-lemma-limit-contained-in-constructible}{{05F4}{source}{}{sourcepointer-2}{}}
\newlabel{limits-lemma-limit-contained-in-constructible@cref}{{[originallemma][0][]05F4}{[1][0][]source}{}{}{}}
\makeatother
\crefname{originallemma}{original lemma}{original lemmas}
\Crefname{originallemma}{Original lemma}{Original lemmas}
\makeatletter
\newlabel{lemma-get-cycles}{{0H59}{source}{}{sourcepointer-3}{}}
\newlabel{lemma-get-cycles@cref}{{[originallemma][0][]0H59}{[1][0][]source}{}{}{}}
\makeatother
\crefname{originalequation}{original equation}{original equations}
\Crefname{originalequation}{Original equation}{Original equations}
\makeatletter
\newlabel{equation-cycle-classes}{{0H58}{source}{}{sourcepointer-4}{}}
\newlabel{equation-cycle-classes@cref}{{[originalequation][0][]0H58}{[1][0][]source}{}{}{}}
\makeatother
\crefname{originalsection}{original section}{original sections}
\Crefname{originalsection}{Original section}{Original sections}
\makeatletter
\newlabel{limits-section-descending-relative}{{01ZL}{source}{}{sourcepointer-5}{}}
\newlabel{limits-section-descending-relative@cref}{{[originalsection][0][]01ZL}{[1][0][]source}{}{}{}}
\makeatother
\crefname{originallemma}{original lemma}{original lemmas}
\Crefname{originallemma}{Original lemma}{Original lemmas}
\makeatletter
\newlabel{limits-lemma-descend-finite-presentation}{{01ZM}{source}{}{sourcepointer-6}{}}
\newlabel{limits-lemma-descend-finite-presentation@cref}{{[originallemma][0][]01ZM}{[1][0][]source}{}{}{}}
\makeatother
\crefname{originallemma}{original lemma}{original lemmas}
\Crefname{originallemma}{Original lemma}{Original lemmas}
\makeatletter
\newlabel{limits-lemma-descend-closed-immersion-finite-presentation}{{01ZP}{source}{}{sourcepointer-7}{}}
\newlabel{limits-lemma-descend-closed-immersion-finite-presentation@cref}{{[originallemma][0][]01ZP}{[1][0][]source}{}{}{}}
\makeatother
\crefname{originallemma}{original lemma}{original lemmas}
\Crefname{originallemma}{Original lemma}{Original lemmas}
\makeatletter
\newlabel{limits-lemma-descend-flat-finite-presentation}{{04AI}{source}{}{sourcepointer-8}{}}
\newlabel{limits-lemma-descend-flat-finite-presentation@cref}{{[originallemma][0][]04AI}{[1][0][]source}{}{}{}}
\makeatother
\crefname{originallemma}{original lemma}{original lemmas}
\Crefname{originallemma}{Original lemma}{Original lemmas}
\makeatletter
\newlabel{limits-lemma-eventually-proper}{{081F}{source}{}{sourcepointer-9}{}}
\newlabel{limits-lemma-eventually-proper@cref}{{[originallemma][0][]081F}{[1][0][]source}{}{}{}}
\makeatother
\crefname{originallemma}{original lemma}{original lemmas}
\Crefname{originallemma}{Original lemma}{Original lemmas}
\makeatletter
\newlabel{limits-lemma-limit-dimension}{{05M5}{source}{}{sourcepointer-10}{}}
\newlabel{limits-lemma-limit-dimension@cref}{{[originallemma][0][]05M5}{[1][0][]source}{}{}{}}
\makeatother
\crefname{originallemma}{original lemma}{original lemmas}
\Crefname{originallemma}{Original lemma}{Original lemmas}
\makeatletter
\newlabel{more-morphisms-lemma-descend-cd}{{0GTM}{source}{}{sourcepointer-11}{}}
\newlabel{more-morphisms-lemma-descend-cd@cref}{{[originallemma][0][]0GTM}{[1][0][]source}{}{}{}}
\makeatother
\crefname{originallemma}{original lemma}{original lemmas}
\Crefname{originallemma}{Original lemma}{Original lemmas}
\makeatletter
\newlabel{lemma-descend-family}{{0H4X}{source}{}{sourcepointer-12}{}}
\newlabel{lemma-descend-family@cref}{{[originallemma][0][]0H4X}{[1][0][]source}{}{}{}}
\makeatother
\crefname{originallemma}{original lemma}{original lemmas}
\Crefname{originallemma}{Original lemma}{Original lemmas}
\makeatletter
\newlabel{lemma-relative-cycles-h-descent}{{0H52}{source}{}{sourcepointer-13}{}}
\newlabel{lemma-relative-cycles-h-descent@cref}{{[originallemma][0][]0H52}{[1][0][]source}{}{}{}}
\makeatother
\title{Relative r-cycles descend through limits of Noetherian schemes}
\author{The Stacks Project Authors}
\date{stacks-project:relative-cycles.tex:local-snapshot-20261012; GNU FDL 1.2 or later}
\begin{document}
\maketitle
\setcounter{subsection}{1}
\noindent\textit{Editorial note.} All delivered mathematical text is the author's own, copied from the pinned Stacks Project chapter file relative-cycles.tex; no statement or proof step was rewritten, reordered or removed. The counted span is the author environment and its own proof environment. References inside the span are handed to the assembler as navigation stubs that carry the canonical tag of the cited tag only; whether each of them is load-bearing is recorded on the stub itself.\par\medskip
\noindent\textit{Allowed standard prerequisites (editorial).}\begin{itemize}
\item The locus where two relative cycles on the fibres agree is closed (relative-cycles chapter, Lemma lemma-relative-cycles-equal, canonical tag 0H5C). cited at proof lines 1398-1400. Two cycles mapping to the same element at the limit base change into its closed locus E in S\_i.
\item Descent of constructible containment along limits: if the limit morphism maps into a constructible set then some transition morphism does (Limits chapter, Lemma limits-lemma-limit-contained-in-constructible, canonical tag 05F4). cited at proof lines 1401-1403. Provides the index j at which the transition morphism lands in E, proving injectivity.
\item Proper completely decomposed covers putting a relative cycle into the image of the cycle-classes map (relative-cycles chapter, Lemma lemma-get-cycles, canonical tag 0H59) and the cycle-classes construction itself (relative-cycles chapter, Equation equation-cycle-classes, canonical tag 0H58). cited at proof lines 1408-1413. Reduces surjectivity to cycles of the form sum n\_a [Z\_a/X'/S']\_r with Z\_a flat of relative dimension at most r.
\item Descending relative algebra through limits: descent of finite presentation (limits-lemma-descend-finite-presentation, 01ZM), of closed immersions of finite presentation (limits-lemma-descend-closed-immersion-finite-presentation, 01ZP), of flatness of finite presentation (limits-lemma-descend-flat-finite-presentation, 04AI), eventual properness (limits-lemma-eventually-proper, 081F), and dimension of the limit (limits-lemma-limit-dimension, 05M5), gathered in Limits, Section limits-section-descending-relative (canonical tag 01ZL). cited at proof lines 1416-1434. Produce an index i with a completely decomposed proper morphism g\_i : S'\_i to S\_i descending g and closed subschemes Z\_\{ai\} flat of relative dimension at most r descending the Z\_a.
\item Descent of completely decomposed morphisms (More on Morphisms, Lemma more-morphisms-lemma-descend-cd, canonical tag 0GTM). cited at proof lines 1433-1434. Makes the descended g\_i completely decomposed.
\item Descent of families of r-cycles along morphisms with separable residue field extensions (relative-cycles chapter, Lemma lemma-descend-family, canonical tag 0H4X). cited at proof lines 1451-1455 and 1458-1459. Its existence and uniqueness produce the unique family alpha\_i with g\_i\textasciicircum{}*alpha\_i = alpha'\_i and force the image of alpha\_i at the limit to be alpha.
\item H-descent for the relative cycle property (relative-cycles chapter, Lemma lemma-relative-cycles-h-descent, canonical tag 0H52). cited at proof lines 1456-1457. Promotes the descended family alpha\_i to an element of z(X\_i/S\_i, r).
\end{itemize}\medskip
\ifdefined\hypertarget\hypertarget{sourcepointer-1}{}\fi
\noindent\textit{Source reference.} Original lemma 0H5C: sources/stacks-chow/tags.txt, source lines 21382--21382. Cited at proof lines 1398-1400: the closed locus E into which the limit morphism maps in the injectivity paragraph.\par\medskip
\ifdefined\hypertarget\hypertarget{sourcepointer-2}{}\fi
\noindent\textit{Source reference.} Original lemma 05F4: sources/stacks-chow/tags.txt, source lines 7013--7013. Cited at proof lines 1401-1403: produces the index j with the transition morphism inside E.\par\medskip
\ifdefined\hypertarget\hypertarget{sourcepointer-3}{}\fi
\noindent\textit{Source reference.} Original lemma 0H59: sources/stacks-chow/tags.txt, source lines 21379--21379. Cited at proof lines 1408-1410: starts the surjectivity paragraph.\par\medskip
\ifdefined\hypertarget\hypertarget{sourcepointer-4}{}\fi
\noindent\textit{Source reference.} Original equation 0H58: sources/stacks-chow/tags.txt, source lines 21378--21378. Cited at proof lines 1410 and 1411: the image of which g\textasciicircum{}*alpha is an element.\par\medskip
\ifdefined\hypertarget\hypertarget{sourcepointer-5}{}\fi
\noindent\textit{Source reference.} Original section 01ZL: sources/stacks-chow/tags.txt, source lines 2598--2598. Cited at proof lines 1416-1417: the machinery used to descend the cover and the flat subschemes.\par\medskip
\ifdefined\hypertarget\hypertarget{sourcepointer-6}{}\fi
\noindent\textit{Source reference.} Original lemma 01ZM: sources/stacks-chow/tags.txt, source lines 2599--2599. Cited at proof lines 1427-1428: descends the morphisms of finite presentation in the diagram.\par\medskip
\ifdefined\hypertarget\hypertarget{sourcepointer-7}{}\fi
\noindent\textit{Source reference.} Original lemma 01ZP: sources/stacks-chow/tags.txt, source lines 2602--2602. Cited at proof lines 1429-1430: descends the closed immersions cutting out the Z\_a.\par\medskip
\ifdefined\hypertarget\hypertarget{sourcepointer-8}{}\fi
\noindent\textit{Source reference.} Original lemma 04AI: sources/stacks-chow/tags.txt, source lines 5617--5617. Cited at proof lines 1430-1431: keeps the descended Z\_\{ai\} flat of relative dimension at most r.\par\medskip
\ifdefined\hypertarget\hypertarget{sourcepointer-9}{}\fi
\noindent\textit{Source reference.} Original lemma 081F: sources/stacks-chow/tags.txt, source lines 10215--10215. Cited at proof lines 1431-1432: descends the properness of the cover g.\par\medskip
\ifdefined\hypertarget\hypertarget{sourcepointer-10}{}\fi
\noindent\textit{Source reference.} Original lemma 05M5: sources/stacks-chow/tags.txt, source lines 7259--7259. Cited at proof lines 1432-1433: controls the relative dimension bound of the descended subschemes.\par\medskip
\ifdefined\hypertarget\hypertarget{sourcepointer-11}{}\fi
\noindent\textit{Source reference.} Original lemma 0GTM: sources/stacks-chow/tags.txt, source lines 20972--20972. Cited at proof lines 1433-1434: makes the descended g\_i completely decomposed.\par\medskip
\ifdefined\hypertarget\hypertarget{sourcepointer-12}{}\fi
\noindent\textit{Source reference.} Original lemma 0H4X: sources/stacks-chow/tags.txt, source lines 21367--21367. Cited at proof lines 1451-1455 and 1458-1459: existence and uniqueness of the descended family alpha\_i.\par\medskip
\ifdefined\hypertarget\hypertarget{sourcepointer-13}{}\fi
\noindent\textit{Source reference.} Original lemma 0H52: sources/stacks-chow/tags.txt, source lines 21372--21372. Cited at proof lines 1456-1457: alpha\_i lies in z(X\_i/S\_i, r).\par\medskip
\begin{lemma}
\label{lemma-descend-through-limit}
Let $S = \lim_{i \in I} S_i$ be the limit of a directed inverse system of
Noetherian schemes with affine transition morphisms.
Let $0 \in I$ and let $X_0 \to S_0$ be a finite type morphism of schemes.
For $i \geq 0$ set $X_i = S_i \times_{S_0} X_0$ and set
$X = S \times_{S_0} X_0$. If $S$ is Noetherian too, then
$$
z(X/S, r) = \colim_{i \geq 0} z(X_i/S_i, r)
$$
where the transition maps are given by base change of relative
$r$-cycles.
\end{lemma}

\begin{proof}
Suppose that $i \geq 0$ and $\alpha_i, \beta_i \in z(X_i/S_i, r)$
map to the same element of $z(X/S, r)$. Then $S \to S_i$
maps into the closed subset $E \subset S_i$ of
Lemma \ref{lemma-relative-cycles-equal}.
Hence for some $j \geq i$ the morphism $S_j \to S_i$
maps into $E$, see Limits, Lemma
\ref{limits-lemma-limit-contained-in-constructible}.
It follows that the base change of
$\alpha_i$ and $\beta_i$ to $S_j$ agree. Thus the map is injective.

\medskip\noindent
Let $\alpha \in z(X/S, r)$. Applying Lemma \ref{lemma-get-cycles}
a completely decomposed proper morphism $g : S' \to S$ such that $g^*\alpha$
is in the image of (\ref{equation-cycle-classes}).
Set $X' = S' \times_S X$. We write $g^*\alpha = \sum n_a [Z_a/X'/S']_r$
for some $Z_a \subset X'$ closed subscheme
flat and of relative dimension $\leq r$ over $S'$.

\medskip\noindent
Now we bring the machinery of
Limits, Section \ref{limits-section-descending-relative} ff to bear.
We can find an $i \geq 0$ such that there exist
\begin{enumerate}
\item a completely decomposed proper morphism $g_i : S'_i \to S_i$
whose base change to $S$ is $g : S' \to S$,
\item setting $X'_i = S'_i \times_{S_i} X_i$
closed subschemes $Z_{ai} \subset X'_i$ flat and of
relative dimension $\leq r$ over $S'_i$ whose base change to $S'$
is $Z_a$.
\end{enumerate}
To do this one uses Limits, Lemmas
\ref{limits-lemma-descend-finite-presentation},
\ref{limits-lemma-descend-closed-immersion-finite-presentation},
\ref{limits-lemma-descend-flat-finite-presentation},
\ref{limits-lemma-eventually-proper}, and
\ref{limits-lemma-limit-dimension}
and More on Morphisms, Lemma
\ref{more-morphisms-lemma-descend-cd}.
Consider
$\alpha'_i = \sum n_a [Z_{ai}/X'_i/S'_i]_r \in z(X'_i/S'_i, r)$.
The image of $\alpha'_i$ in $z(X'/S', r)$ agrees with the base change
$g^*\alpha$ by construction.

\medskip\noindent
Set $S''_i = S'_i \times_{S_i} S'_i$ and $X''_i = S''_i \times_{S_i} X_i$
and set $S'' = S' \times_S S'$ and $X'' = S'' \times_S X$.
We denote $\text{pr}_1, \text{pr}_2 : S'' \to S'$ and
$\text{pr}_1, \text{pr}_2 : S''_i \to S'_i$ the projections.
The two base changes $\text{pr}_1^*\alpha'_i$ and $\text{pr}_1^*\alpha'_i$
map to the same element of $z(X''/S'', r)$ because
$\text{pr}_1^*g^*\alpha = \text{pr}_1^*g^*\alpha$.
Hence after increasing $i$ we may assume that
$\text{pr}_1^*\alpha'_i = \text{pr}_1^*\alpha'_i$
by the first paragraph of the proof.
By Lemma \ref{lemma-descend-family}
we obtain a unique family $\alpha_i$
of $r$-cycles on fibres of $X_i/S_i$
with $g_i^*\alpha_i = \alpha'_i$ (this uses that $S'_i \to S_i$
is completely decomposed).
By Lemma \ref{lemma-relative-cycles-h-descent}
we see that $\alpha_i \in z(X_i/S_i, r)$.
The uniqueness in Lemma \ref{lemma-descend-family} implies that the
image of $\alpha_i$ in $z(X/S, r)$ is $\alpha$ and the proof is complete.
\end{proof}

\end{document}
